\section{Background and Related Work}
\label{sec:background}

\subsection{Leader-Based and Leaderless Consensus}

Leader-based protocols such as Raft~\cite{ongaro2014raft} designate one
replica as the leader: client requests are directed to it, it appends
operations to its log and replicates them to followers, and an entry is
committed once a quorum has acknowledged it. This centralises the write path
and makes leader failure require an election before processing resumes.
Paxos~\cite{lamport1998parttime} established the foundational agreement
framework, but its abstract presentation leaves substantial implementation
decisions to system designers, so systems in the same protocol family can
differ considerably in execution path, storage, and communication behaviour.

Leaderless protocols such as EPaxos~\cite{moraru2013epaxos} let any replica
propose operations and determine dependencies among them. Commands that do
not conflict can commit without a single global ordering point, at the cost of
dependency tracking and conflict detection: replicas must still agree on
dependencies and execute conflicting operations in a compatible order. The
loss of one replica does not require electing a replacement before the others
continue proposing.

\subsection{Related Work}

Comparative studies of consensus designs have reached different conclusions
depending on the systems and workloads examined. Howard and Mortier found
that the choice of implementation and configuration often matters more than
the protocol family itself when comparing Paxos and Raft
\cite{howard2020raft}. van Renesse et al. showed that the relative
performance of Paxos and leaderless replication depends on the workload,
particularly on command conflict \cite{vanrenesse2015paxos}. Tollman et al.
re-evaluated EPaxos under geo-replicated workloads and found that the
original benchmarking approach did not expose the performance impact of
conflicting commands, with substantially worse tail latency than the
reported median \cite{tollman2021epaxos}. These results motivate evaluating
multiple independent implementations of each protocol family under
controlled conditions, as we do here.

Our study compares two Raft implementations (HashiCorp Raft and etcd Raft)
with two EPaxos implementations (the original \texttt{efficient/epaxos} and
\texttt{nvanbenschoten/epaxos}) while varying replica count, concurrency,
communication delay, and failure conditions, so the results describe concrete
implementations under a controlled workload rather than protocol-independent
properties of Raft or EPaxos.
